% 87-02 ①(87쪽) — 보기 그래프: $x=-1$ 에서 극대, $x=3$ 에서 기울기 0 인 위로 볼록한 사차함수 개형. 그림 자체가 보기라 TikZ 로 다시 그림.
% 라벨은 원문 그대로 -1, O, 3, x, y, y=f(x) 와 점선 두 줄만.
\begin{tikzpicture}[baseline=(current bounding box.center), x=0.36cm, y=0.36cm, line width=0.5pt]
  \path (-2.9,-2.4) rectangle (5.15,3.5);
  \draw[->, >=stealth, line width=0.7pt] (-2.8,0) -- (4.85,0) node[below=1pt, font=\footnotesize] {$x$};
  \draw[->, >=stealth, line width=0.7pt] (0,-2.1) -- (0,3.2) node[left=1pt, font=\footnotesize] {$y$};
  \draw[densely dotted, line width=0.4pt] (-1,0) -- (-1,2.4);
  \draw[densely dotted, line width=0.4pt] (3,0) -- (3,0.35);
  \draw[line width=0.6pt, smooth] plot coordinates {(-2.6,-1.6) (-2.45,-0.9) (-2.3,-0.3) (-2.1,0.45) (-1.85,1.2) (-1.6,1.8) (-1.35,2.2) (-1.15,2.38) (-1,2.4) (-0.8,2.35) (-0.5,2.2) (-0.1,1.95) (0.4,1.6) (0.9,1.25) (1.4,0.95) (1.9,0.7) (2.4,0.5) (2.8,0.38) (3,0.35) (3.25,0.3) (3.45,0.18) (3.6,-0.05) (3.8,-0.5) (3.95,-1.1) (4.05,-1.7)};
  \node[below=1pt, font=\footnotesize] at (-1.3,0) {$-1$};
  \node[below left=-2pt and -3pt, font=\footnotesize] at (0,0) {$\pt{O}$};
  \node[below=1pt, font=\footnotesize] at (3,0) {$3$};
  \node[anchor=west, font=\footnotesize] at (0.25,2.95) {$y=f(x)$};
\end{tikzpicture}
